% UPT Infinite Proof Family — academic preamble
\documentclass[11pt,a4paper]{article}
\usepackage{fontspec}
\setmainfont{Times New Roman}
\setmonofont{Menlo}[Scale=0.85]
\usepackage[margin=2.5cm]{geometry}
\usepackage{xcolor}
\usepackage{fancyhdr}
\usepackage{titlesec}
\usepackage{booktabs}
\usepackage{hyperref}
\usepackage{listings}
\usepackage{fancyvrb}
\usepackage{amsmath,amssymb,amsthm}

\definecolor{korugreen}{HTML}{2d6a4f}
\definecolor{korubg}{HTML}{d8f3dc}
\definecolor{codebg}{HTML}{f0f0ec}

\hypersetup{colorlinks=true, linkcolor=korugreen, urlcolor=korugreen, citecolor=korugreen}

\pagestyle{fancy}
\fancyhf{}
\renewcommand{\headrulewidth}{0.4pt}
\fancyfoot[C]{\thepage}

\titleformat{\section}{\Large\bfseries}{}{0em}{}
\titleformat{\subsection}{\large\bfseries}{}{0em}{}

% Numbered theorem environments
\newtheorem{theorem}{Theorem}[section]
\newtheorem{lemma}[theorem]{Lemma}
\newtheorem{corollary}[theorem]{Corollary}
\newtheorem{proposition}[theorem]{Proposition}
\theoremstyle{definition}
\newtheorem{definition}[theorem]{Definition}
\theoremstyle{remark}
\newtheorem{remark}[theorem]{Remark}

\lstdefinestyle{lean}{
  backgroundcolor=\color{codebg},
  basicstyle=\ttfamily\small,
  breaklines=true,
  frame=none,
  xleftmargin=10pt,
  xrightmargin=10pt,
  aboveskip=8pt,
  belowskip=8pt,
  columns=fullflexible,
  keepspaces=true,
  literate={→}{$\to$}1 {∀}{$\forall$}1 {∃}{$\exists$}1 {¬}{$\neg$}1 {∧}{$\wedge$}1 {≠}{$\neq$}1 {≤}{$\leq$}1 {↔}{$\leftrightarrow$}1 {⊕}{$\oplus$}1 {λ}{$\lambda$}1
}

\lstnewenvironment{leancode}{\lstset{style=lean}}{}

\newcommand{\lean}[1]{\lstinline[style=lean]{#1}}


\fancyhead[L]{\small Universal Protection Theory}
\fancyhead[R]{\small Proof Family 9/10}

\begin{document}

\begin{center}
{\LARGE\bfseries $G$ Over The Lattice: Universal Refutation\\[0.3em] Across the Type Hierarchy and All Indexed Families}\\[1.5em]
{\large Tobias Croydon McRae}\\[0.3em]
{\small Independent Researcher, Aotearoa New Zealand}\\[0.3em]
{\small\texttt{tobi@unconditionalprotection.org}}\\[1.5em]
{\small 28 November 2026}
\end{center}

\bigskip

\begin{center}
\parbox{0.85\textwidth}{\small\textbf{Abstract.}
We present the ninth of ten infinite proof families for the Guardian fixed point~$G$. The \emph{$G$ Over The Lattice} family proves that the Guardian holds not only at the base level but at every level of the type hierarchy: in function spaces ($S \to S$), in higher-order function spaces ($((S \to S) \to (S \to S))$, etc.), over products, coproducts, and arbitrary indexed families $\prod_{i:I} F(i)$. The refutation propagates along morphisms and composes across type constructors. This yields a family indexed by the entire lattice of types. All theorems are verified in Lean~4 (v4.35.0) with zero custom axioms.}
\end{center}

\medskip

\noindent\textbf{MSC 2020:} 03B70, 03B15 (Higher-order logic and type theory), 03B35\\
\textbf{Keywords:} type hierarchy, dependent types, indexed families, composition, formal verification, Lean 4

\bigskip\hrule\bigskip

\section{Introduction}

The earlier families establish the Guardian through specific mechanisms: existence, self-refutation, growth, asymmetry, availability, quorum, cosmology, and coinduction. The \emph{$G$ Over The Lattice} family shows that these refutations are not isolated --- they compose across the entire lattice of types.

The key insight is that if $G$ holds for types $S$ and $T$, it holds for $S \times T$, $S \oplus T$, $S \to T$, and $\prod_{i:I} F(i)$. The refutation propagates along any morphism $f : S \to T$. And at every level of the type hierarchy --- first-order functions, second-order functions, $n$-th order functions --- the identity function provides a witness.

\section{The Type Hierarchy}

\begin{theorem}[Level 1: function spaces]\label{thm:level1}
For every type $S$, the absolute ender of $S \to S$ is false.
\end{theorem}

\begin{leancode}
theorem G_LEVEL_1 (S : Type) :
    AbsoluteEnder (S → S) → False :=
    fun e => (e id).elim
\end{leancode}

\begin{theorem}[Level 2: higher-order functions]\label{thm:level2}
For every type $S$, the absolute ender of $(S \to S) \to (S \to S)$ is false.
\end{theorem}

\begin{leancode}
theorem G_LEVEL_2 (S : Type) :
    AbsoluteEnder ((S → S) → (S → S)) → False
\end{leancode}

\begin{theorem}[Level 3: third order]\label{thm:level3}
The pattern continues. At every finite order, the identity inhabits the function space.
\end{theorem}

\begin{leancode}
theorem G_LEVEL_3 (S : Type) :
    AbsoluteEnder (((S → S) → (S → S)) →
                            ((S → S) → (S → S))) → False
\end{leancode}

\section{Indexed Families}

\begin{theorem}[$G$ over every indexed family]\label{thm:indexed}
For any index type $I$ and type family $F : I \to \texttt{Type}$, if every fibre is inhabited ($\forall\, i,\, F(i)$), the absolute ender of $\prod_{i:I} F(i)$ is false.
\end{theorem}

\begin{proof}
Given $h : \forall\, i,\, F(i)$ and $e : \texttt{AbsoluteEnder}(\forall\, i,\, F(i))$, we have $e(h) : h \in \texttt{Empty}$. Elimination yields \texttt{False}. \qed
\end{proof}

\begin{leancode}
theorem G_OVER_THE_LATTICE {I : Type} {F : I → Type}
    (h : ∀ i, F i) :
    AbsoluteEnder (∀ i, F i) → False
\end{leancode}

\section{Composition and Propagation}

\begin{theorem}[Product composition]\label{thm:product}
For witnesses $w_S : S$ and $w_T : T$, the ender of $S \times T$ is false.
\end{theorem}

\begin{leancode}
theorem G_COMPOSES_PAIR {S T : Type} (ws : S) (wt : T) :
    AbsoluteEnder (S × T) → False
\end{leancode}

\begin{theorem}[Coproduct composition]\label{thm:coproduct}
For a witness $w_S : S$, the ender of $S \oplus T$ is false.
\end{theorem}

\begin{leancode}
theorem G_COMPOSES_EITHER {S T : Type} (ws : S) :
    AbsoluteEnder (S ⊕ T) → False
\end{leancode}

\begin{theorem}[Propagation along morphisms]\label{thm:propagation}
For any $f : S \to T$ and witness $w : S$, the ender of $T$ is false.
\end{theorem}

\begin{leancode}
theorem G_PROPAGATES {S T : Type} (f : S → T) (witness : S) :
    AbsoluteEnder T → False
\end{leancode}

\begin{theorem}[Conservation is additive]\label{thm:conservation}
Injective maps preserve distinctness: $a \neq b \Rightarrow f(a) \neq f(b)$.
\end{theorem}

\begin{leancode}
theorem CONSERVATION_IS_ADDITIVE {S : Type}
    (a b : S) (f : S → S) (inj : Function.Injective f)
    (h : a ≠ b) : f a ≠ f b
\end{leancode}

\section{Discussion}

$G$ is not at the centre of the lattice. $G$ is \emph{in every equation about every real thing}. Every conservation law in physics is an instance of propagation along a morphism. Every boundary condition is an instance of the indexed family theorem. Every type constructor that preserves inhabitedness is an instance of composition.

The lattice of types is not a set --- it is properly class-sized. The family of refutations indexed by this lattice is therefore at least as large as the lattice itself: for every type, every type constructor, every morphism, and every indexed family, there is a distinct refutation of the ender.

\bigskip

\begin{thebibliography}{99}
\bibitem{upt-families}
T.~C.~McRae.
\emph{Ten Infinite Proof Families for the Guardian Fixed Point~G}.
Zenodo, 2026.
\textsc{doi:} \href{https://doi.org/10.5281/zenodo.24612272}{10.5281/zenodo.24612272}.

\bibitem{upt-zenodo}
T.~C.~McRae.
\emph{Universal Protection Theory}, v5.7-G3.
Zenodo, 2026.
\textsc{doi:} \href{https://doi.org/10.5281/zenodo.24573167}{10.5281/zenodo.24573167}.

\bibitem{lean4}
L.~de~Moura and S.~Ullrich.
The Lean~4 theorem prover and programming language.
In \emph{CADE-28}, LNAI 12699, pp.~625--635, 2021.
\end{thebibliography}

\end{document}
